% Chapter 28: Language Workshop: One Passage, Every Skill

\chapter{Language Workshop: One Passage, Every Skill}

\section{Pedagogical Purpose}
This chapter takes integration one step further: instead of starting from an error-filled passage, it starts from a single short reading passage and uses that one passage as the shared source for every kind of task in this book: finding vocabulary in context, identifying grammar at work, transforming sentences, editing, and writing a response. This mirrors how a "language workshop" period works in a real classroom.

\begin{learningobjectives}
The learner can:
\begin{itemize}
    \item read a short passage and answer simple comprehension questions about it.
    \item find a synonym, antonym, or word meaning using context clues from the passage.
    \item identify a named grammar feature (a tense, a type of sentence, a part of speech) within the passage.
    \item transform a sentence taken from the passage without changing its meaning.
    \item edit a short passage using the five-step checklist from Chapter 27.
    \item write a short piece that responds to or is modelled on the workshop passage.
\end{itemize}
\end{learningobjectives}

\section{Prerequisite Check}
\begin{quickcheck}
Look back at Chapter 27's five-step checklist. Without looking, try to name as many of the five steps as you can.
\end{quickcheck}

\begin{rememberbox}
You already have every tool this chapter needs: vocabulary skills from Chapters 17-19, grammar knowledge from every earlier chapter, the editing checklist from Chapter 27, and paragraph-writing skills from Chapter 24. This chapter simply brings them all to the same short passage, one task at a time.
\end{rememberbox}

\section{Chapter Opener}
Ms. Sharma has photocopied a short passage for the whole class and pinned extra questions around the room — some about meaning, some about grammar, some about writing.

\textbf{Ms. Sharma:} "Today we are not learning anything brand new. Instead, we are going to treat this one short passage like a workshop table. We'll pick up a different tool each time — first a vocabulary tool, then a grammar tool, then an editing tool, then a writing tool — but we'll keep coming back to the very same passage."

\textbf{Rohan:} "So it's like using the same photo for lots of different questions?"

\textbf{Maya:} "Yes — and I bet the passage will start to feel really familiar by the end!"

\textbf{Ms. Sharma:} "Exactly, Maya. Let's begin."

\section{Language Experience}
\textbf{The Workshop Passage: "The Kite Festival"}

\begin{quote}
Every January, the rooftops of Rohan's town fill with colour. Families climb up before sunrise, carrying kites of every shape and size. Rohan's grandfather, who has flown kites since he was a boy, taught him how to release the string at just the right moment. "Feel the wind first," his grandfather always says, "before you let go." By afternoon, the sky is crowded with paper diamonds, dipping and soaring above the narrow streets. Rohan's kite, a bright red one shaped like a fish, climbed higher than any other kite on the roof. He grinned with pride, and for a moment, he forgot to hold on tightly — the string slipped, and his kite drifted silently away over the rooftops.
\end{quote}
\begin{tcolorbox}[colback=gray!5, colframe=gray!50, sharp corners, boxsep=0pt, left=5pt, right=5pt, top=5pt, bottom=5pt, breakable]
Every January, the rooftops of Rohan's town fill with colour. Families climb up before sunrise, carrying kites of every shape and size. Rohan's grandfather, who has flown kites since he was a boy, taught him how to release the string at just the right moment. "Feel the wind first," his grandfather always says, "before you let go." By afternoon, the sky is crowded with paper diamonds, dipping and soaring above the narrow streets. Rohan's kite, a bright red one shaped like a fish, climbed higher than any other kite on the roof. He grinned with pride, and for a moment, he forgot to hold on tightly --- the string slipped, and his kite drifted silently away over the rooftops.
\end{tcolorbox}

\section{Look and Notice}
\begin{thinkbox}
\begin{enumerate}
    \item Who are the characters in this passage, and where does it take place?
    \item Find one word in the passage you think another word could replace without changing the meaning much.
    \item Find one sentence written as a piece of direct speech. How do you know it is direct speech?
    \item Find one verb in the past tense and one verb in the present tense. Why do you think the passage uses both?
\end{enumerate}
\end{thinkbox}

\section{Discover / Explicit Teaching}
A language workshop moves through the same passage several times, each time with a different purpose:

\begin{table}[h!]
\centering
\begin{tabular}{|l|p{5cm}|p{4cm}|}
\hline
\textbf{Pass} & \textbf{Purpose} & \textbf{Skill Area} \\
\hline
1 & Read for meaning & Comprehension \\
2 & Notice vocabulary and context clues & Vocabulary (Ch. 17-19) \\
3 & Identify grammar features at work & Grammar (whole book) \\
4 & Transform sentences from the passage & Sentence craft (Ch. 20-22) \\
5 & Edit a related passage & Editing (Ch. 27) \\
6 & Write a response & Writing (Ch. 24, 26) \\
\hline
\end{tabular}
\end{table}

Rereading the same text for different purposes is a genuine skill — a single passage almost always has more to notice than a single reading reveals.

\section{Grammar Word}
\begin{grammarword}{CONTEXT CLUE}
A \textbf{context clue} is a hint in the surrounding sentence or passage that helps a reader work out the meaning of an unfamiliar word, without needing a dictionary.

\textit{Examples:}
\begin{itemize}
    \item "The sky was \textbf{crowded} with paper diamonds, dipping and soaring above the streets." — the surrounding words (many kites, a busy sky) suggest "crowded" means "full of many things."
    \item "Rohan's kite \textbf{drifted} silently away." — "silently" and "away" suggest slow, gentle movement, helping explain "drifted."
\end{itemize}

\textit{Quick Check:} In the workshop passage, what context clues help you guess the meaning of "soaring"?
\end{grammarword}

\section{Core Explanation}
\subsection*{Reading for vocabulary in context}
Example: "Families climb up before sunrise" — the surrounding words (rooftops, kites, January) help confirm "climb" means going upward. \textit{Check:} Can I guess a sensible meaning for an unfamiliar word just from the words around it?

\subsection*{Reading for grammar features}
Example: "Rohan's kite... \textbf{climbed} higher" is past tense; "his grandfather always \textbf{says}" is present tense. \textit{Check:} Can I name the tense, sentence type, or part of speech being asked about?

\subsection*{Transforming sentences from a real passage}
Example: Active — "His grandfather taught him..." $\rightarrow$ Passive — "He was taught by his grandfather..." \textit{Check:} Does my transformed sentence still mean exactly what the original passage sentence meant?

\subsection*{Editing within a workshop task}
A related passage is deliberately seeded with errors for practice. \textit{Check:} Am I applying the full five-step checklist from Chapter 27?

\subsection*{Writing a response to the passage}
Example: Continuing "The Kite Festival" by imagining where the lost kite lands. \textit{Check:} Does my writing genuinely respond to or extend the passage?

\section{Worked Example}
\begin{examplebox}
\textbf{Task:} Using the workshop passage, (1) explain the meaning of "grinned" using context, and (2) transform the sentence "His grandfather taught him how to release the string" into the passive voice.

\textit{Step 1 (Vocabulary):} Look at the surrounding words — "He grinned with pride." The context (pride, success) suggests "grinned" means smiled widely, happily.
\textit{Step 2 (Grammar transformation):} Identify the doer ("his grandfather") and receiver ("him").
\textit{Step 3:} Make the receiver the new subject and add the correct form of "be" plus the past participle.

\textbf{Answer:}
1. "Grinned" means smiled widely and happily, shown by the context clue "with pride."
2. "He was taught by his grandfather how to release the string."
\end{examplebox}

\section{Guided Practice}
\begin{activitybox}{Activity A: Identification}
Using the workshop passage, find one example of each.
\begin{enumerate}
    \item A word that could be replaced by a synonym without changing the meaning much.
    \item A verb in the simple present tense.
    \item A piece of direct speech.
\end{enumerate}
\end{activitybox}

\section{Independent Practice}
\begin{activitybox}{Independent Practice}
\begin{enumerate}
    \item Find two more words in the workshop passage and explain their meaning using a context clue.
    \item Imagine Rohan's grandfather also said, "I am proud of you." Change this into reported speech.
    \item Identify one adjective and one adverb from the passage, and explain what each one describes.
    \item \textit{Reasoning:} Why might rereading the same passage several times help you notice more than reading it only once?
\end{enumerate}
\end{activitybox}

\section{Summary}
\begin{summarybox}
\begin{itemize}
    \item A \textbf{language workshop} returns to the same short passage several times, each time with a different purpose.
    \item \textbf{Context clues} are hints in the surrounding text that help explain an unfamiliar word's meaning.
    \item \textbf{Comprehension questions} ask about meaning; \textbf{language-focus questions} ask about specific grammar or vocabulary.
    \item Sentences from a passage can be \textbf{transformed} without changing their meaning.
    \item Writing tasks should genuinely \textbf{respond to or extend} the passage.
\end{itemize}
\end{summarybox}

\section{Key Terms}
\begin{table}[h!]
\centering
\begin{tabular}{|l|p{9cm}|}
\hline
\textbf{Term} & \textbf{Simple Meaning} \\
\hline
Language workshop & Using one passage for several different language tasks \\
Context clue & A hint in the surrounding text that helps explain a word's meaning \\
Comprehension question & A question about what a passage means \\
Language-focus question & A question about a specific grammar or vocabulary feature \\
Transformation & Rewriting a sentence in a different grammatical form without changing its meaning \\
\hline
\end{tabular}
\end{table}

\begin{selfcheck}
I can:
\begin{itemize}
    \item read a short passage and answer comprehension questions about it.
    \item use context clues to explain an unfamiliar word's meaning.
    \item identify a named grammar feature within a passage.
    \item transform a sentence from a passage without changing its meaning.
    \item edit a related passage using the five-step checklist.
    \item write a short piece that responds to or extends a passage.
\end{itemize}
\end{selfcheck}